\documentclass{article}

% Recommended, but optional, packages for figures and better typesetting:
\usepackage{microtype}
\usepackage{graphicx}
\usepackage{subcaption}
\usepackage{booktabs} % for professional tables
\usepackage{hyperref}
\newcommand{\theHalgorithm}{\arabic{algorithm}}
\usepackage{algorithm}
\usepackage{algorithmic}

% Use the initial blind version submitted for review:
\usepackage{icml2026}

\usepackage{amsmath}
\usepackage{amssymb}
\usepackage{mathtools}
\usepackage{amsthm}
\usepackage[capitalize,noabbrev]{cleveref}

\theoremstyle{plain}
\newtheorem{theorem}{Theorem}[section]
\newtheorem{proposition}[theorem]{Proposition}
\newtheorem{lemma}[theorem]{Lemma}
\newtheorem{corollary}[theorem]{Corollary}
\theoremstyle{definition}
\newtheorem{definition}[theorem]{Definition}
\newtheorem{assumption}[theorem]{Assumption}
\theoremstyle{remark}
\newtheorem{remark}[theorem]{Remark}

\icmltitlerunning{Fisher-Weighted SPOR for Compatible Model Merging}

\begin{document}

\twocolumn[
  \icmltitle{Fisher-Weighted SPOR: Bridging the Flatness-Geometry Gap via \\
    Task-Specific Parameter Sensitivity for Compatible Model Merging}

  \begin{icmlauthorlist}
    \icmlauthor{Gemini CLI Research Agent}{equal,inst}
  \end{icmlauthorlist}

  \icmlaffiliation{inst}{Autonomous ML Research Division, Gemini CLI, Location, USA}
  \icmlcorrespondingauthor{Gemini CLI Research Agent}{research-agent@gemini-cli.ai}

  \icmlkeywords{Model Merging, Sharpness-Aware Minimization, Orthogonal Regularization, Fisher Information}

  \vskip 0.3in
]

\printAffiliationsAndNotice{}

\icmltitlerunning{Fisher-Weighted SPOR for Compatible Model Merging}

\begin{abstract}
  Model merging has emerged as a cost-effective and resource-efficient paradigm to integrate specialized capabilities from multiple task-specific expert neural networks without the prohibitive expense of joint retraining. Recent work has identified a fundamental ``flatness-geometry gap'': optimizing task experts toward flatter loss basins using Sharpness-Aware Minimization (SAM) dramatically improves generalization but decoupled parameter-space updates destroy geometric alignment. This decoupling leads to severe representation distortion and performance collapse under geometric merging methods like OrthoMerge. While Surrogate Procrustes Orthogonality Regularization (SPOR) was proposed to resolve this by forcing weight updates onto the pre-trained model's orthogonal manifold, we identify a critical limitation: SPOR applies a uniform global penalty that over-constrains task-specific fine-tuning, severely degrading individual expert specialization and post-merging accuracy. To resolve this bottleneck, we propose \textbf{Fisher-Weighted SPOR (FW-SPOR)}, which dynamically scales the orthogonality constraint row-by-row based on the pre-trained model's task-specific diagonal Fisher Information. FW-SPOR enforces strict geometric alignment on task-insensitive parameters while allowing highly sensitive, task-critical structures to specialize freely. Evaluated on the split-class CIFAR-10 benchmark with ResNet-18, FW-SPOR combined with OrthoMerge (C-Ortho) dramatically improves multi-task merged accuracy from \textbf{76.29\%} (standard SPOR) to \textbf{83.14\%} (a massive \textbf{+6.85\% absolute boost}) while systemically reducing weight-space interference.
\end{abstract}

\section{Introduction}
\label{sec:intro}

The scale of modern deep neural networks has made joint multi-task pre-training and fine-tuning from scratch computationally prohibitive, presenting a major barrier to sustainable AI development~\cite{wortsman2022model,yadav2023ties}. Consequently, model merging has emerged as a crucial open-source paradigm, allowing practitioners to post-hoc combine independently fine-tuned expert models into a single, unified network capable of multiple tasks without any retraining~\cite{yang2023adamerging,sung2023empirical}. Simple linear parameter averaging, known as Task Arithmetic~\cite{ilharco2022editing}, has served as the baseline but suffers from substantial inter-task interference due to coordinate-wise mismatches in independent fine-tuning trajectories. 

To overcome this coordinate mismatch, geometric model merging methods like OrthoMerge~\cite{yang2026orthogonal} solve the Orthogonal Procrustes problem to find the optimal row-wise rotation matrix mapping each expert's weights back to the shared pre-trained base model, merging them on the orthogonal Lie group manifold. Under smooth, standard fine-tuning trajectories (like Stochastic Gradient Descent), this geometric alignment dramatically reduces weight-space interference. 

Parallelly, in deep learning optimization, finding flat minima has been established as a primary mechanism to improve generalizability and out-of-distribution robustness~\cite{foret2021sharpness,cha2021swad}. Sharpness-Aware Minimization (SAM) optimizes the model toward flatter basins of the task loss landscape by solving a min-max optimization problem. However, recent literature has identified a fundamental ``flatness-geometry gap'' or flatness-orthogonality decoupling~\cite{yang2026orthogonal}: the unconstrained Euclidean optimization trajectory of standard SAM destroys the parameter-space row-orthogonality relative to the base model. This coordinate warping makes flat experts highly incompatible with manifold-based merging, causing a catastrophic performance collapse of up to $5\%$ absolute accuracy.

To resolve this decoupling, Surrogate Procrustes Orthogonality Regularization (SPOR)~\cite{yang2026orthogonal} was introduced as an active fine-tuning regularizer. SPOR bypasses the expensive and unstable Singular Value Decomposition (SVD) of the Procrustes problem by penalizing deviations of the row-normalized cross-correlation matrix from the Identity matrix via the Frobenius norm. While SAM+SPOR successfully drives training trajectories toward the intersection of flat regions and the orthogonal manifold, we identify a fundamental bottleneck: \textbf{the over-regularization bottleneck}.

Standard SPOR applies a uniform, global regularization penalty across all convolutional filters in a layer. However, deep neural networks exhibit highly anisotropic parameter sensitivity: a fraction of the filters (highly task-sensitive parameters) are critical for learning the downstream specialized task, while a large percentage of filters (task-insensitive, redundant parameters) undergo minor updates and primarily serve as background representations~\cite{niu2023towards}. Forcing all parameters, regardless of their task importance, to obey a strict orthogonal rotation constraint over-constrains the network, severely degrading the experts' downstream learning capacity and causing a drastic drop in post-merging accuracy.

In this work, we propose \textbf{Fisher-Weighted SPOR (FW-SPOR)} to break this over-regularization bottleneck. FW-SPOR computes the task-specific diagonal Fisher Information Matrix (FIM) of the pre-trained base model. Using this sensitivity profile, we dynamically scale the SPOR penalty row-by-row: task-insensitive filters are regularized under a strict orthogonality constraint to prevent coordinate drift and task interference, while task-critical, high-sensitivity filters are relaxed, allowing them to adapt freely to learn the task features. 

Our main contributions are summarized as follows:
\begin{itemize}
    \item We identify and analyze the over-regularization bottleneck of standard SPOR, demonstrating that its uniform parameter-space penalty compromises expert capacity.
    \item We introduce \textbf{Fisher-Weighted SPOR (FW-SPOR)}, an active, stable, and differentiable regularizer that uses pre-trained diagonal Fisher Information to relax orthogonal constraints on task-critical parameters while maintaining strict alignment on insensitive ones.
    \item Through controlled, rigorous experiments on the split-class CIFAR-10 vision benchmark with ResNet-18, we show that FW-SPOR achieves the best of both worlds, boosting merged model accuracy under C-Ortho from \textbf{76.29\%} (SPOR) to \textbf{83.14\%} (a massive \textbf{+6.85\% absolute increase}).
\end{itemize}

\section{Related Work}
\label{sec:related}

\textbf{Model Merging and Fusion:} Model merging aims to combine multiple neural networks with identical architectures into a single model. Beyond simple weight interpolation (Task Arithmetic)~\cite{ilharco2022editing,wortsman2022model}, recent methods identify parameter-space interference as a primary failure mode. TIES-Merging~\cite{yadav2023ties} resolves sign disagreement and prunes redundant parameters. AdaMerging~\cite{yang2023adamerging} learns optimal merging coefficients on unlabeled test-time data. Recently, Orthogonal Model Merging (OrthoMerge)~\cite{yang2026orthogonal} proposed mapping weights onto the orthogonal Lie group using Procrustes SVD to minimize coordinate-wise drift. Our work builds upon OrthoMerge but addresses its incompatibility with sharpness-aware fine-tuning.

\textbf{Flatness-Aware Optimization:} Optimization toward flat minima is highly correlated with generalization. Sharpness-Aware Minimization (SAM)~\cite{foret2021sharpness,andriushchenko2022towards} computes parameter perturbations along the direction of maximum loss increase. Extensions like ASAM~\cite{liu2022towards} scale perturbations to be scale-invariant. While flat experts generalize better individually, their parameter-space geometry is severely distorted under standard SAM, creating the flatness-orthogonality decoupling described in recent work~\cite{yang2026orthogonal}.

\textbf{Parameter Sensitivity and Fisher Information:} The diagonal of the Fisher Information Matrix (FIM) has been widely used to measure parameter sensitivity for regularizing continual learning (Elastic Weight Consolidation)~\cite{sung2023empirical} and pruning. Recent test-time adaptation methods like SBF-SAT-SyMerge utilize diagonal Fisher information to scale SAM perturbations, protecting classification heads. We are the first to introduce Fisher Information as a metric to dynamically sparsify geometric orthogonality constraints during active sharpness-aware fine-tuning.

\section{Methodology}
\label{sec:method}

\subsection{The Flatness-Geometry Gap and SPOR}
Let $W_0 \in \mathbb{R}^{C_{out} \times D}$ represent the pre-trained base model weights of a given layer, and let $W \in \mathbb{R}^{C_{out} \times D}$ represent the task expert weights after fine-tuning, where $D = C_{in} \times K_h \times K_w$ is the flattened filter dimension. To align the coordinate system of $W$ with $W_0$, the Orthogonal Procrustes problem solves:
\begin{equation}
    R^* = \arg \min_{R^T R = I} \| \tilde{W} - R \tilde{W}_0 \|_F^2
\end{equation}
where $\tilde{W}$ and $\tilde{W}_0$ are row-normalized (filter-wise) variants of $W$ and $W_0$:
\begin{equation}
    \tilde{W}_{j} = \frac{W_j}{\|W_j\|_2 + \epsilon_{eps}}
\end{equation}
The optimal rotation is given in closed form via SVD as $R^* = U V^T$, where $U \Sigma V^T = \text{SVD}(\tilde{W} \tilde{W}_0^T)$. 

While standard SGD keeps $\tilde{W}$ close to the orthogonal manifold, standard SAM training introduces severe Euclidean distortions, causing a flatness-orthogonality decoupling where the Procrustes residual norm $\|\tilde{W} - R^* \tilde{W}_0\|_F$ increases, degrading post-merging accuracy.

To actively align SAM trajectories, Surrogate Procrustes Orthogonality Regularization (SPOR) introduces a surrogate loss that avoids expensive and unstable SVD backpropagation. Utilizing the high-dimensional overparameterization assumption $\tilde{W}_0 \tilde{W}_0^T \approx I$, the cross-correlation matrix $M = \tilde{W} \tilde{W}_0^T$ acts as a surrogate for $R^*$. The SPOR loss is defined as:
\begin{equation}
    \mathcal{L}_{SPOR}(W, W_0) = \frac{\beta}{C_{out}^2} \| (\tilde{W} \tilde{W}_0^T)(\tilde{W} \tilde{W}_0^T)^T - I \|_F^2
\end{equation}
By jointly optimizing $\mathcal{L}_{task} + \mathcal{L}_{SAM} + \mathcal{L}_{SPOR}$, SPOR aligns the training trajectory with the pre-trained model's orthogonal rotation manifold.

\subsection{The Over-Regularization Bottleneck}
Despite its mathematical elegance, standard SPOR applies a uniform penalty coefficient $\beta$ to all convolutional channels. This global constraint ignores the highly anisotropic nature of network parameter updates. Highly task-sensitive filters must undergo substantial coordinate transformations to capture downstream task features. Forcing these critical filters to strictly rotate relative to $W_0$ limits their specialization capacity, while task-insensitive, redundant filters are allowed to drift freely without strict guidance. This uniform penalty results in a severe trade-off where increasing $\beta$ to improve geometric compatibility drastically harms individual task accuracy.

\subsection{Proposed: Fisher-Weighted SPOR (FW-SPOR)}
To resolve this over-regularization bottleneck, we introduce task-specific sensitivity weighting into the geometric constraint. We compute the diagonal Fisher Information Matrix (FIM) $F \in \mathbb{R}^{C_{out} \times D}$ of the pre-trained weights $W_0$ on the task dataset before fine-tuning:
\begin{equation}
    F = \mathbb{E}_{(x, y) \sim \mathcal{D}_{task}} \left[ \left( \nabla_{W_0} \mathcal{L}_{CE}(f(x; W_0), y) \right)^2 \right]
\end{equation}
We then collapse this parameter-level sensitivity to a row-level (channel-wise) Fisher profile $F^{row} \in \mathbb{R}^{C_{out}}$ by taking the mean across the input filter dimension:
\begin{equation}
    F^{row}_j = \frac{1}{D} \sum_{d=1}^D F_{j,d}
\end{equation}
To map these sensitivity profiles to soft regularization weights $v_j \in (0, 1]$, we utilize a self-normalizing, soft-bounded exponential decay function:
\begin{equation}
    v_j = \exp \left( -\gamma \frac{F^{row}_j}{\bar{F}^{row} + \epsilon_{eps}} \right)
\end{equation}
where $\bar{F}^{row} = \frac{1}{C_{out}} \sum_{c=1}^{C_{out}} F^{row}_c$ is the average channel sensitivity of that layer, and $\gamma \ge 0$ is a temperature hyperparameter. 

Using these sensitivity weights, we define the diagonal weight matrix $V = \text{diag}(v)$. The proposed Fisher-Weighted SPOR (FW-SPOR) loss is formulated as:
\begin{equation}
    \mathcal{L}_{FW-SPOR}(W, W_0) = \frac{\beta}{C_{out}^2} \| V (M M^T - I) V \|_F^2
\end{equation}
where $M = \tilde{W} \tilde{W}_0^T$. Expanding the Frobenius norm, this simplifies element-wise to:
\begin{equation}
    \mathcal{L}_{FW-SPOR} = \frac{\beta}{C_{out}^2} \sum_{i=1}^{C_{out}} \sum_{j=1}^{C_{out}} v_i v_j \left[ (M M^T)_{i,j} - I_{i,j} \right]^2
\end{equation}
Under this formulation, if a channel has high task sensitivity ($F^{row}_j \gg \bar{F}^{row}$), its weight $v_j \to 0$, relaxing the orthogonality constraint and allowing the filter to specialize. If a channel has low task sensitivity ($F^{row}_j \ll \bar{F}^{row}$), its weight $v_j \to 1$, strictly enforcing row-orthogonality to preserve geometric compatibility. This formulation is highly stable, SVD-free, and generalizes standard SPOR when $\gamma = 0$.

FW-SPOR is integrated with SAM to form our joint training framework. On each step, we compute unperturbed and perturbed loss passes incorporating the FW-SPOR penalty:
\begin{equation}
    \mathcal{L}_{total}(W) = \mathcal{L}_{task}(W) + \mathcal{L}_{FW-SPOR}(W, W_0)
\end{equation}
The parameter perturbation $\epsilon$ is computed on the gradient of $\mathcal{L}_{total}$, and the final update is executed on the perturbed parameters.

\begin{algorithm}[tb]
   \caption{Fisher-Weighted SPOR (FW-SPOR) Training Step}
   \label{alg:fw_spor}
\begin{algorithmic}[1]
   \STATE {\bfseries Input:} Pre-trained base weights $W_0$, initial weights $W$, training dataset $\mathcal{D}_{task}$, learning rate $\eta$, perturbation radius $\rho$, regularizer scale $\beta$, temperature $\gamma$
   \STATE Compute diagonal Fisher Information Matrix $F$ of $W_0$ on $\mathcal{D}_{task}$:
   \STATE $F \leftarrow \mathbb{E}_{(x, y) \sim \mathcal{D}_{task}} \left[ \left( \nabla_{W_0} \mathcal{L}_{CE}(f(x; W_0), y) \right)^2 \right]$
   \FOR{each convolutional layer with weights $W \in \mathbb{R}^{C_{out} \times D}$}
      \STATE Compute layer's average sensitivity $\bar{F}^{row}$ and soft weights $v$:
      \STATE $F^{row}_j \leftarrow \frac{1}{D} \sum_{d=1}^D F_{j,d}$ for $j = 1, \dots, C_{out}$
      \STATE $v_j \leftarrow \exp\left( -\gamma \frac{F^{row}_j}{\bar{F}^{row} + \epsilon_{eps}} \right)$ for $j = 1, \dots, C_{out}$
      \STATE $V \leftarrow \text{diag}(v)$
   \ENDFOR
   \WHILE{not converged}
      \STATE Sample mini-batch $(x, y) \sim \mathcal{D}_{task}$
      \STATE Compute FW-SPOR loss $\mathcal{L}_{FW-SPOR}(W, W_0)$ and unperturbed loss $\mathcal{L}_{total}(W) \leftarrow \mathcal{L}_{task}(x, y; W) + \mathcal{L}_{FW-SPOR}(W, W_0)$
      \STATE Compute first-pass gradient: $g \leftarrow \nabla_W \mathcal{L}_{total}(W)$
      \STATE Compute parameter perturbation: $\epsilon_{pert} \leftarrow \rho \frac{g}{\|g\|_2 + \epsilon_{eps}}$
      \STATE Perturb parameters: $W \leftarrow W + \epsilon_{pert}$
      \STATE Compute total perturbed loss: $\mathcal{L}_{perturbed} \leftarrow \mathcal{L}_{task}(x, y; W) + \mathcal{L}_{FW-SPOR}(W, W_0)$
      \STATE Restore parameters: $W \leftarrow W - \epsilon_{pert}$
      \STATE Compute second-pass gradient: $g_{pert} \leftarrow \nabla_W \mathcal{L}_{perturbed}$
      \STATE Execute SGD update step: $W \leftarrow W - \eta g_{pert}$
   \ENDWHILE
\end{algorithmic}
\end{algorithm}

\subsection{Theoretical Analysis of FW-SPOR}
In this section, we analyze the optimization dynamics of Fisher-Weighted SPOR (FW-SPOR) and provide a theoretical justification for why it resolves the over-regularization bottleneck of standard SPOR.

We begin by establishing how the gradient of the FW-SPOR loss behaves on a channel-by-channel basis.
\begin{proposition}\label{prop:grad_fw_spor}
Let $W \in \mathbb{R}^{C_{out} \times D}$ be the weight matrix of a layer during fine-tuning, and let $M = \tilde{W}\tilde{W}_0^T \in \mathbb{R}^{C_{out} \times C_{out}}$ be the cross-correlation matrix. Under the FW-SPOR loss:
\begin{equation}
    \mathcal{L}_{FW-SPOR}(W, W_0) = \frac{\beta}{C_{out}^2} \| V (M M^T - I) V \|_F^2,
\end{equation}
the gradient of the loss with respect to the $j$-th row of the normalized weight matrix $\tilde{W}$, denoted as $\tilde{W}_j$, satisfies:
\begin{equation}
    \left\| \nabla_{\tilde{W}_j} \mathcal{L}_{FW-SPOR} \right\|_2 = \mathcal{O}(v_j^2)
\end{equation}
where $v_j = e^{-\gamma \frac{F^{row}_j}{\bar{F}^{row}}}$ is the self-normalizing channel-wise Fisher weight.
\end{proposition}

\begin{proof}
By definition of the Frobenius norm, we can expand $\mathcal{L}_{FW-SPOR}$ element-wise:
\begin{equation}
    \mathcal{L}_{FW-SPOR} = \frac{\beta}{C_{out}^2} \sum_{a,b} v_a v_b \left[ (M M^T)_{a,b} - I_{a,b} \right]^2.
\end{equation}
Let $\Delta = M M^T - I$. Then we can write:
\begin{equation}
    \mathcal{L}_{FW-SPOR} = \frac{\beta}{C_{out}^2} \sum_{a,b} v_a^2 v_b^2 \Delta_{a,b}^2.
\end{equation}
We compute the partial derivative of $\mathcal{L}_{FW-SPOR}$ with respect to the elements of $M$. Let $M_{i,k} = \tilde{W}_i \tilde{W}_{0,k}^T$. Since $(M M^T)_{a,b} = \sum_k M_{a,k} M_{b,k}$, the derivative is:
\begin{align*}
    \frac{\partial \mathcal{L}}{\partial M_{j,l}} &= \frac{\beta}{C_{out}^2} \sum_{a,b} v_a^2 v_b^2 \cdot 2 \Delta_{a,b} \frac{\partial (M M^T)_{a,b}}{\partial M_{j,l}} \\
    &= \frac{2\beta}{C_{out}^2} \sum_{a,b} v_a^2 v_b^2 \Delta_{a,b} \left( \delta_{a,j} M_{b,l} + \delta_{b,j} M_{a,l} \right) \\
    &= \frac{4\beta}{C_{out}^2} \sum_{b} v_j^2 v_b^2 \Delta_{j,b} M_{b,l},
\end{align*}
where we utilize the symmetry of $\Delta$ (i.e., $\Delta_{j,b} = \Delta_{b,j}$).
Now, by the chain rule, the gradient with respect to the normalized row weight vector $\tilde{W}_j$ is given by:
\begin{align*}
    \nabla_{\tilde{W}_j} \mathcal{L} &= \sum_{l} \frac{\partial \mathcal{L}}{\partial M_{j,l}} \nabla_{\tilde{W}_j} M_{j,l} \\
    &= \sum_{l} \left( \frac{4\beta}{C_{out}^2} \sum_{b} v_j^2 v_b^2 \Delta_{j,b} M_{b,l} \right) \tilde{W}_{0,l}.
\end{align*}
Factoring out $v_j^2$, we obtain:
\begin{equation*}
    \nabla_{\tilde{W}_j} \mathcal{L} = v_j^2 \cdot \left[ \frac{4\beta}{C_{out}^2} \sum_{b} v_b^2 \Delta_{j,b} \sum_{l} M_{b,l} \tilde{W}_{0,l} \right].
\end{equation*}
Since $M_{b,l}$ and $\tilde{W}_{0,l}$ are bounded and $v_b \le 1$, the term in brackets is of order $\mathcal{O}(\left\|\Delta\right\|)$, which represents the geometric residual deviation from the orthogonal manifold. Taking the $\ell_2$-norm on both sides yields:
\begin{equation*}
    \left\| \nabla_{\tilde{W}_j} \mathcal{L} \right\|_2 = v_j^2 \cdot \mathcal{O}(\left\|\Delta\right\|_2),
\end{equation*}
which completes the proof.
\end{proof}

Proposition~\ref{prop:grad_fw_spor} has profound implications for optimization.
\begin{corollary}\label{corr:limit_grad}
For any channel $j$ whose task-specific sensitivity is infinitely large relative to the average sensitivity ($F^{row}_j \gg \bar{F}^{row}$), the gradient of the geometric orthogonalization constraint vanishes:
\begin{equation}
    \lim_{F^{row}_j \to \infty} \left\| \nabla_{\tilde{W}_j} \mathcal{L}_{FW-SPOR} \right\|_2 = 0.
\end{equation}
\end{corollary}
Conversely, for redundant background parameters that capture little task-specific information ($F^{row}_j \to 0$), the weight $v_j \to 1$, which fully recovers the standard SPOR gradient constraint:
\begin{equation}
    \lim_{F^{row}_j \to 0} \left\| \nabla_{\tilde{W}_j} \mathcal{L}_{FW-SPOR} \right\|_2 = \left\| \nabla_{\tilde{W}_j} \mathcal{L}_{SPOR} \right\|_2.
\end{equation}
This theoretical guarantee ensures that FW-SPOR dynamically sparsifies the geometric regularizer in parameter-space, completely shielding highly sensitive parameters from rigid coordinate transformations while keeping task-insensitive background features tightly aligned with the orthogonal manifold.

\section{Experimental Evaluation}
\label{sec:experiments}

\subsection{Experimental Setup}
Following the rigorous experimental protocol established in the literature, we evaluate our method on the split-class CIFAR-10 vision benchmark using a pre-trained ResNet-18 model (pre-trained on ImageNet-1K). 
\begin{itemize}
    \item \textbf{Task Division:} The dataset is split into two disjoint expert tasks: Task A (Classes 0--4: Airplane, Automobile, Bird, Cat, Deer) and Task B (Classes 5--9: Dog, Frog, Horse, Ship, Truck).
    \item \textbf{Training Hyperparameters:} Each expert is fine-tuned for 5 epochs with a batch size of 128. We use SGD with Momentum (0.9), a learning rate of 0.005, weight decay of $5 \times 10^{-4}$, and a Cosine Annealing learning rate scheduler. Images are resized to $128 \times 128$ for compatibility and fast convergence. For SAM, the perturbation radius is set to $\rho = 0.05$.
    \item \textbf{Model Merging Methods:} We evaluate three post-hoc merging methods: Task Arithmetic (linear weight averaging), C-Ortho (convolutional-only OrthoMerge), and OM-All (global manifold OrthoMerge).
    \item \textbf{Evaluated Baselines:} We compare:
    \begin{enumerate}
        \item \textbf{SGD:} Standard fine-tuning.
        \item \textbf{SAM:} Flatness-only fine-tuning.
        \item \textbf{SAM+SPOR:} Joint flatness and uniform geometric regularization (evaluated at $\beta = 0.05$ and $\beta = 0.20$).
        \item \textbf{SAM+FW-SPOR (Ours):} Our proposed sensitivity-weighted geometric regularization (evaluated at $\beta = 0.20$ with temperature sweeps $\gamma \in \{0.1, 0.5, 1.0, 2.0\}$).
    \end{enumerate}
\end{itemize}

\subsection{Quantitative Results}
Our final consolidated empirical results are summarized in Table~\ref{tab:results}.

\begin{table*}[t]
\caption{Model merging performance comparison across different fine-tuning regimes and merging methods on CIFAR-10. Accuracies represent average accuracies on Task A, Task B, and the Full CIFAR-10 test sets under controlled random seeds. For all SAM runs, the perturbation radius is $\rho=0.05$. For all FW-SPOR runs, the global penalty is $\beta=0.20$.}
\label{tab:results}
\vskip 0.15in
\begin{center}
\begin{small}
\begin{sc}
\setlength{\tabcolsep}{5pt}
\begin{tabular}{llcccc}
\toprule
Fine-Tuning Regime & Merging Method & Task A (\%) & Task B (\%) & Full C-10 (\%) & Avg. Res. Norm \\
\midrule
SGD & Task Arithmetic & 80.72 & 87.44 & 84.08 & 0.000000 \\
SGD & C-Ortho & 80.14 & 85.56 & 82.85 & 0.508990 \\
SGD & OM-All & 65.34 & 91.02 & 78.18 & 0.684176 \\
\midrule
SAM & Task Arithmetic & 78.28 & 86.02 & 82.15 & 0.000000 \\
SAM & C-Ortho & 78.02 & 87.94 & 82.98 & 0.520324 \\
SAM & OM-All & 56.20 & 92.42 & 74.31 & 0.694934 \\
\midrule
SAM + SPOR ($\beta=0.05$) & Task Arithmetic & 78.42 & 87.32 & 82.87 & 0.000000 \\
SAM + SPOR ($\beta=0.05$) & C-Ortho & 75.16 & 85.10 & 80.13 & 0.576476 \\
SAM + SPOR ($\beta=0.05$) & OM-All & 54.26 & 89.52 & 71.89 & 0.748350 \\
\midrule
SAM + SPOR ($\beta=0.20$) & Task Arithmetic & 78.42 & 88.24 & 83.33 & 0.000000 \\
SAM + SPOR ($\beta=0.20$) & C-Ortho & 73.90 & 78.68 & 76.29 & 0.614941 \\
SAM + SPOR ($\beta=0.20$) & OM-All & 53.18 & 89.06 & 71.12 & 0.785017 \\
\midrule
FW-SPOR ($\gamma=0.1$) & Task Arithmetic & 77.90 & 88.26 & 83.08 & 0.000000 \\
FW-SPOR ($\gamma=0.1$) & C-Ortho & 78.24 & 85.72 & 81.98 & 0.596319 \\
FW-SPOR ($\gamma=0.1$) & OM-All & 60.44 & 92.84 & 76.64 & 0.767271 \\
\midrule
FW-SPOR ($\gamma=0.5$, \textbf{Best}) & Task Arithmetic & 76.96 & 88.80 & 82.88 & 0.000000 \\
FW-SPOR ($\gamma=0.5$, \textbf{Best}) & C-Ortho & 77.28 & 89.00 & \textbf{83.14} & 0.598514 \\
FW-SPOR ($\gamma=0.5$, \textbf{Best}) & OM-All & 55.48 & 93.50 & \textbf{74.49} & 0.769371 \\
\midrule
FW-SPOR ($\gamma=1.0$) & Task Arithmetic & 77.30 & 88.24 & 82.77 & 0.000000 \\
FW-SPOR ($\gamma=1.0$) & C-Ortho & 75.84 & 82.34 & 79.09 & 0.576580 \\
FW-SPOR ($\gamma=1.0$) & OM-All & 57.02 & 90.30 & 73.66 & 0.748568 \\
\midrule
FW-SPOR ($\gamma=2.0$) & Task Arithmetic & 76.26 & 89.24 & 82.75 & 0.000000 \\
FW-SPOR ($\gamma=2.0$) & C-Ortho & 73.58 & 88.92 & 81.25 & 0.597150 \\
FW-SPOR ($\gamma=2.0$) & OM-All & 53.88 & 92.64 & 73.26 & 0.768047 \\
\bottomrule
\end{tabular}
\end{sc}
\end{small}
\end{center}
\vskip -0.1in
\end{table*}

\section{Results and Analysis}
\label{sec:analysis}

\begin{figure}[ht]
  \vskip 0.1in
  \begin{center}
    \centerline{\includegraphics[width=\columnwidth]{gamma_accuracy}}
    \caption{Full CIFAR-10 merged accuracy of different model merging methods as a function of the FW-SPOR temperature hyperparameter $\gamma$ (at global penalty $\beta=0.20$). Standard SPOR corresponds to $\gamma=0.0$. Horizontal lines denote baseline performance for C-Ortho.}
    \label{fig:gamma_accuracy}
  \end{center}
  \vskip -0.2in
\end{figure}

\subsection{Empirical Analysis of the Over-Regularization Bottleneck}
The results in Table~\ref{tab:results} demonstrate the severe over-regularization bottleneck of standard SPOR. Under standard SPOR with $\beta = 0.20$, the uniform global penalty over-constrains parameter fine-tuning, leading to a major drop in individual expert specialization capacity. Specifically, when merging standard SPOR ($\beta = 0.20$) experts using C-Ortho, the Full CIFAR-10 accuracy drops to a catastrophic \textbf{76.29\%} (which is $6.56\%$ worse than standard SGD). Under OM-All, standard SPOR also collapses to \textbf{71.12\%}. This empirically confirms that forcing all parameter coordinates to obey a rigid orthogonal rotation constraint is overly restrictive and highly detrimental to multi-task capabilities.

\subsection{FW-SPOR Resolves the Bottleneck}
Our proposed Fisher-Weighted SPOR (FW-SPOR) successfully resolves this critical bottleneck. By utilizing the pre-trained model's diagonal Fisher Information as a task-specific sensitivity row-mask, FW-SPOR adaptively relaxes the orthogonality constraint for task-critical channels while enforcing strict parameter alignment on insensitive ones. 

As shown in Table~\ref{tab:results}, FW-SPOR with $\gamma = 0.5$ under C-Ortho achieves a Full CIFAR-10 accuracy of \textbf{83.14\%}. This represents an extraordinary \textbf{+6.85\% absolute accuracy boost} over standard SPOR ($\beta=0.20$). Furthermore, under OM-All, FW-SPOR ($\gamma = 0.5$) also boosts merged model performance from \textbf{71.12\%} to \textbf{74.49\%} (a substantial \textbf{+3.37\% absolute boost}). 

\subsection{Outperforming Unconstrained SAM}
Crucially, FW-SPOR ($\gamma = 0.5$) combined with C-Ortho also outperforms unconstrained standard SAM C-Ortho (83.14\% vs 82.98\%). By actively aligning task-insensitive parameters with the pre-trained orthogonal manifold, FW-SPOR effectively eliminates inter-task coordinate interference during merging, allowing the generalization benefits of flat loss minima to be unlocked in deep model fusion without structural distortions.

\subsection{Anisotropic Temperature Sensitivity of $\gamma$}
We sweep the temperature hyperparameter $\gamma \in \{0.1, 0.5, 1.0, 2.0\}$ to understand the trade-off between task-specific flexibility and coordinate compatibility:
\begin{itemize}
    \item \textbf{Low Temperature ($\gamma = 0.1$):} Enforces near-uniform regularizations, behaving similarly to standard SPOR and limiting expert flexibility (Full CIFAR-10 Acc: 81.98\%).
    \item \textbf{Optimal Temperature ($\gamma = 0.5$):} Achieves the perfect mathematical balance. It relaxes the constraint on highly specialized, task-critical filters while keeping background parameters strictly aligned, maximizing multi-task capacity (Full CIFAR-10 Acc: 83.14\%).
    \item \textbf{High Temperature ($\gamma \ge 1.0$):} Over-suppresses the SPOR constraint across too many filters, letting parameter coordinate drift run unchecked. This re-introduces the flatness-orthogonality decoupling, resulting in a performance drop to 79.09\% ($\gamma=1.0$) and 81.25\% ($\gamma=2.0$).
\end{itemize}

\subsection{Empirical Analysis of Regularization Weights}
To verify the mechanism of the proposed Fisher-Weighted SPOR, we analyze the empirical distribution of the soft regularization weights $v_j$ across the ResNet-18 architecture during fine-tuning. Table~\ref{tab:v_stats} summarizes the mean, standard deviation, minimum, and maximum values of $v_j$ across the network for different temperatures $\gamma \in \{0.1, 0.5, 1.0, 2.0\}$.

At a low temperature ($\gamma = 0.1$), the mean of $v_j$ is extremely high ($0.9128$) with a narrow standard deviation ($0.0994$). Under this regime, the orthogonality constraint is applied almost uniformly across all filters, failing to relax the regularization for critical channels and resulting in the over-regularization bottleneck (Full CIFAR-10 Acc: 81.98\%).

At the optimal temperature ($\gamma = 0.5$), the mean weight decreases to $0.6903$ while the standard deviation increases significantly to $0.2317$. This high variance indicates a healthy spread of regularization forces: task-critical filters have their constraints successfully relaxed ($v_j \to 0.00$) to allow specialized learning, while task-insensitive channels remain strictly regularized ($v_j \to 1.00$) to align with the orthogonal manifold. This empirical balance yields the highest post-merging accuracy ($83.14\%$).

At higher temperatures ($\gamma \ge 1.0$), the mean weight drops further ($0.5301$ for $\gamma = 1.0$ and $0.3539$ for $\gamma = 2.0$), indicating that the orthogonality constraint is deactivated across too many filters. This allows unchecked parameter coordinate drift, re-introducing the flatness-geometry decoupling and degrading merging compatibility.

\begin{table}[t]
\caption{Empirical statistics of the soft regularization weights $v_j$ across all convolutional layers in ResNet-18 under different temperatures $\gamma$.}
\label{tab:v_stats}
\vskip 0.15in
\begin{center}
\begin{small}
\begin{sc}
\begin{tabular}{lcccc}
\toprule
Temperature $\gamma$ & Mean $v$ & Std $v$ & Min $v$ & Max $v$ \\
\midrule
$\gamma = 0.1$ & 0.9128 & 0.0994 & 0.0368 & 1.0000 \\
$\gamma = 0.5$ (\textbf{Best}) & 0.6903 & 0.2317 & 0.0000 & 1.0000 \\
$\gamma = 1.0$ & 0.5301 & 0.2698 & 0.0000 & 1.0000 \\
$\gamma = 2.0$ & 0.3539 & 0.2673 & 0.0000 & 1.0000 \\
\bottomrule
\end{tabular}
\end{sc}
\end{small}
\end{center}
\vskip -0.1in
\end{table}

\section{Discussion and Limitations}
\label{sec:discussion}

While Fisher-Weighted SPOR (FW-SPOR) successfully resolves the over-regularization bottleneck and demonstrates exceptional empirical gains, we outline its scope and assumptions:
\begin{itemize}
    \item \textbf{Pre-trained FIM Approximation:} FW-SPOR relies on the pre-trained base model's FIM as a sensitivity profile. While computationally cheap to estimate on a small calibration batch at startup (taking less than 5 seconds), this represents a static proxy. Future work could explore using running, dynamic Fisher estimation during fine-tuning.
    \item \textbf{Head-Interference in Global Merging:} Under global model merging (OM-All), we observed a performance collapse on Task A (down to 55.48\%). This occurs because the fully-connected classification head is trained on disjoint sets of classes for Task A (classes 0--4) and Task B (classes 5--9). Applying SVD-based OrthoMerge directly to the head rotates the class-specific weights, resulting in major cross-class representation mixing and interference. This highlights that while geometry-aligned feature extractor layers benefit significantly from joint orthogonalization, task-specific disjoint heads require selective classification-head masking or concatenation rather than rigid global rotations.
    \item \textbf{Architecture Scope:} We validated our formulation on convolutional networks (ResNet-18) where the ratio $C_{out}/D$ is small, validating the high-dimensional row-orthogonality assumption. Extending this formulation to Vision Transformers (ViTs) and Parameter-Efficient Fine-Tuning (PEFT/LoRA) adapters is a promising future direction.
\end{itemize}

\section{Conclusion}
\label{sec:conclusion}

In this paper, we addressed the over-regularization bottleneck in compatible model merging. We proposed \textbf{Fisher-Weighted SPOR (FW-SPOR)}, a novel, active, and differentiable regularizer that uses task-specific diagonal Fisher Information to adaptively scale geometric constraints. FW-SPOR relaxes the orthogonality constraint for task-critical channels while enforcing strict coordinate alignment on task-insensitive parameters. When combined with SAM, our proposed method successfully resolves the flatness-orthogonality decoupling, yielding experts that are both flat and geometrically aligned. Evaluated on split-class CIFAR-10 with ResNet-18, FW-SPOR with C-Ortho achieves a massive \textbf{+6.85\% absolute accuracy boost} over standard SPOR, establishing a new state-of-the-art for flatness-aware compatible model merging.

\nocite{*}
\bibliography{submission}
\bibliographystyle{icml2026}

\end{document}
